% Fragmento público generado desde propietarios íntegros.
% Residencia: 52.2.1.
% Fuente: ../incoming/radion_monografia_20260827/manuscrito/sections/03_cierre_exacto.tex:1-50
\noindent El cierre exacto se obtiene mediante dos secciones construidas antes del cociclo que las relaciona.\par\medskip

\subsubsection{Construcción independiente de las secciones directa y torsional}

\paragraph{Sección directa del sistema temporal dodecafásico}

La rueda dodecafásica tiene centros
\begin{equation}
\theta_m=30m-10^\circ,\qquad m=1,\ldots,12.
\label{eq:rueda-dodecafasica}
\end{equation}
La segunda fase es \(\theta_2=50^\circ\). Junto a la publicación
\(A=10^3\alphah\), forma la sección directa
\begin{align}
S_E
&=\frac{T_+}{360}\,(50+10^3\alphah)\\
&=2\pih\left(\frac5{36}+\frac{25}{9}\alphah\right).
\label{eq:seccion-electrica}
\end{align}
El subíndice \(E\) significa aquí \emph{directa o común}; su realización
eléctrica llegará después del lector constitutivo. No se postula
literalmente \(A=E\).

Los cilindros coinductivos prueban
\begin{equation}
1<S_E<2.
\label{eq:SE-intervalo}
\end{equation}
La división APP fuerza entonces
\begin{equation}
q_E=\lfloor S_E\rfloor=1,
\qquad
D_E=\operatorname{rem}_1(S_E)=S_E-1.
\label{eq:DE}
\end{equation}
El entero uno no es el número de vueltas de la hélice. Es el cociente único
de la sección directa dentro de la celda de unidad.

\paragraph{Sección torsional determinada por el operador \(\Delta_4\)}

La segunda sección utiliza una salida coinductiva anterior:
\begin{equation}
\Delta_4=
\frac{\pih-e_{\HMT}\log\pih}{270}
\left(1+\frac{\pih}{729}\right).
\label{eq:delta4}
\end{equation}
La palabra \emph{torsión} designa aquí el cuanto global de la carta HMT. No
lo identifica automáticamente con la torsión de Cartan
\(T^I=D_\omega e^I\). El puente hacia Cartan exige otra flecha tipada.
